% part: intuitionistic-logic
% chapter: propositions-as-types
% section: normalization

\documentclass[../../../include/open-logic-section]{subfiles}

\begin{document}

\olfileid{int}{pty}{nor}

\olsection{স্বাভাবিকীকরণ}

এই অংশে প্রমাণ করব, হ্রাসের কোনো একটি ক্রমে প্রত্যেক
নিষ্পাদনকে স্বাভাবিক নিষ্পাদনে হ্রাস করা যায়। একেই
\emph{স্বাভাবিকীকরণ ধর্ম} বলে। সূত্র-হিসেবে-টাইপ অনুরূপতা
ব্যবহার করব: প্রত্যেক প্রমাণপদকে স্বাভাবিক রূপে হ্রাস করা
যায় দেখালে স্বাভাবিক নিষ্পাদনের !!{derivation}s-এর
স্বাভাবিকীকরণও পাওয়া যাবে।

প্রথমে পদগুলির জটিলতা মাপার কয়েকটি অপেক্ষক সংজ্ঞায়িত করি।
!!a{formula}s-এর \emph{দৈর্ঘ্য}~$\len{!A}$ হলো
\begin{align*}
  \len{p} &= 0\\
  \len{!A \land !B} &= \len{!A} + \len{!B} + 1\\
  \len{!A \lor !B} &= \len{!A} + \len{!B} + 1 \\
  \len{!A \lif !B} &= \len{!A} + \len{!B} + 1.
\end{align*}
রিডেক্স (হ্রাসযোগ্য পদাংশ)~$M$-এর জটিলতা তার
\emph{কর্তনমাত্রা}~$\cutrank{M}$ দিয়ে মাপি:
\begin{align*}
  \cutrank{(\lambd[\typeof{x}{!A}][\typeof{N}{!B}])Q} &=
  \len{!A} + \len{!B} + 1 \\
  \cutrank{\ande{i}{\andi{\typeof{M}{!A}}{\typeof{N}{!B}}}} &=
  \len{!A} + \len{!B} + 1 \\
  \cutrank{\ore{\ori{i}{!A_{i}}{\typeof{M}{!A_i}}}
    {\typeof{x_1}{!A_1}}{\typeof{N_1}{!C}}
    {\typeof{x_2}{!A_2}}{\typeof{N_2}{!C}}} &=
  \len{!A_1} + \len{!A_2} + 1
\end{align*}
% BN-SRC-604 TARGET-CORRECTION-BEGIN
\par\smallskip\noindent\textnormal{\small\textbf{উৎসসংশোধন (BN-SRC-604):} শেষ সারির সমষ্টি-টাইপের দুই শাখা \(!A_1\) ও \(!A_2\); উৎসে অনির্ধারিত \(!A\) ও \(!B\) ছিল।}
% BN-SRC-604 TARGET-CORRECTION-END
প্রমাণপদের জটিলতা তার মধ্যে সবচেয়ে জটিল
রিডেক্স দিয়ে মাপি; পদটি স্বাভাবিক হলে মান~$0$:
\[
  \maxrank{M} = \max \{\cutrank{N} | N \text{ হলো } M\text{-এর উপপদ ও রিডেক্স}\}
\]
% BN-NORM-177 TARGET-NORMALIZATION-BEGIN
\par\smallskip\noindent\textnormal{\small\textbf{ভাষাগত স্বাভাবিকীকরণ (BN-NORM-177):} সর্বাধিক কর্তনমাত্রার সেট-সংজ্ঞার ইংরেজি শর্ত বাংলা করা হয়েছে; গাণিতিক শর্ত অপরিবর্তিত।}
% BN-NORM-177 TARGET-NORMALIZATION-END

\begin{lem}\ollabel{lem:subst}
  $\Subst{M}{\typeof{N}{!A}}{\typeof{x}{!A}}$ রিডেক্স
  হলে এবং $M \not\ident x$ হলে নিচের
  অবস্থাগুলির একটি ঘটে:
  \begin{enumerate}
  \item $M$ নিজেই রিডেক্স; অথবা
  \item $M$-এর রূপ $\ande{i}{x}$, আর $N$-এর রূপ
    $\andi{P_1}{P_2}$
  \item $M$-এর রূপ $\ore{i}{x_1}{P_1}{x_2}{P_2}$, আর $N$-এর
    রূপ $\ori{i}{}{Q}$
  \item $M$-এর রূপ $x Q$, আর $N$-এর রূপ
    $\lambd[x][P]$
  \end{enumerate}
  প্রথম অবস্থায় $\cutrank{\Subst{M}{N}{x}} = \cutrank{M}$;
  অন্য অবস্থাগুলিতে $\cutrank{\Subst{M}{N}{x}} = \len{!A}$।
% BN-SRC-605 TARGET-CORRECTION-BEGIN
\par\smallskip\noindent\textnormal{\small\textbf{উৎসসংশোধন (BN-SRC-605):} শেষ সমীকরণে উৎসের অতিরিক্ত বন্ধনী বাদ দেওয়া হয়েছে।}
% BN-SRC-605 TARGET-CORRECTION-END
\end{lem}
\begin{proof}
  $M$-এর গঠন অনুযায়ী আরোহের সাহায্যে প্রমাণ।
  \begin{enumerate}
  \item $M$ একক চলরাশি $y$ হলে এবং $y \not\ident x$ হলে
    $\Subst{y}{N}{x}$ হলো $y$; তাই সেটি রিডেক্স নয়।
  \item $M$-এর রূপ $\andi{N_1}{N_2}$,
    $\lambd[x][N]$ অথবা $\ori{i}{!A}{N}$ হলে
    $\Subst{M}{\typeof{N}{!A}}{\typeof{x}{!A}}$-ও
    সেই রূপের, তাই রিডেক্স নয়।
  \item $M$-এর রূপ $\ande{i}{P}$ হলে দুটি অবস্থা দেখি।
    \begin{enumerate}
      \item $P$-এর রূপ $\andi{P_1}{P_2}$ হলে $M \ident
        \ande{i}{\andi{P_1}{P_2}}$ রিডেক্স; স্পষ্টতই
        \[
        \Subst{M}{N}{x} \ident
        \ande{i}{\andi{\Subst{P_1}{N}{x}}{\Subst{P_2}{N}{x}}}
        \]
        রিডেক্স। দুটির কর্তনমাত্রা সমান।
      \item $P$ একক চলরাশি হলে প্রতিস্থাপনের ফল
        রিডেক্স হতে গেলে সেটি $x$ হতে হবে,
        আর $N$-এর রূপ $\andi{P_1}{P_2}$ হতে হবে। এবার
        \[
        \Subst{M}{N}{x} \ident \Subst{\ande{i}{x}}{\andi{P_1}{P_2}}{x},
        \]
        সেটি $\ande{i}{\andi{P_1}{P_2}}$। এর কর্তনমাত্রা
        $\len{!A}$।
% BN-SRC-606 TARGET-CORRECTION-BEGIN
\par\smallskip\noindent\textnormal{\small\textbf{উৎসসংশোধন (BN-SRC-606):} চলরাশি \(x\) রিডেক্স নয়; তার অনির্ধারিত কর্তনমাত্রার উল্লেখ সরিয়ে \(x\)-এর টাইপের দৈর্ঘ্য রাখা হয়েছে।}
% BN-SRC-606 TARGET-CORRECTION-END
    \end{enumerate}
  \end{enumerate}
  $\ore{N}{x_1}{N_1}{x_2}{N_2}$ ও $P Q$-এর
  ক্ষেত্রগুলিও অনুরূপ।
\end{proof}

\begin{lem}
  $M$ সংকুচিত হয়ে $M'$ হলে, এবং $M$-এর প্রত্যেক
  প্রকৃত রিডেক্স উপপদ~$N$-এর জন্য
  $\cutrank{M} > \cutrank{N}$ হলে, $\cutrank{M} >
  \maxrank{M'}$.
\end{lem}

\begin{proof}
  অবস্থাভেদে প্রমাণ।
  \begin{enumerate}
  \item $M$-এর রূপ $\ande{i}{\andi{M_1}{M_2}}$ হলে
    $M'$ হলো $M_i$। $M_i$-এর প্রত্যেক উপপদই
    $M$-এর প্রকৃত উপপদ, তাই দাবিটি সিদ্ধ।
  \item $M$-এর রূপ
    $(\lambd[\typeof{x}{!A}][N])\typeof{Q}{!A}$ হলে $M'$ হলো
    $\Subst{N}{\typeof{Q}{!A}}{\typeof{x}{!A}}$।
    $M'$-এর কোনো রিডেক্স নিই। হয় $N$-এ
    তার অনুরূপ সমান কর্তনমাত্রার পদাংশ আছে, যা
    অনুমান অনুসারে $\cutrank{M}$-এর চেয়ে কম;
    নয়তো কর্তনমাত্রা $\len{!A}$, যা সংজ্ঞা অনুযায়ী
    $\cutrank{(\lambd[x^{!A}][N])Q}$-এর চেয়ে কম।
  \item $M$-এর রূপ
    \[
    \ore{\ori{i}{}{\typeof{N}{!A_i}}}
        {\typeof{x_1}{!A_1}}{\typeof{N_1}{!C}}
        {\typeof{x_2}{!A_2}}{\typeof{N_2}{!C}},
    \]
    হলে $M' \ident \Subst{N_i}{N}{\typeof{x_i}{!A_i}}$।
    $M'$-এর কোনো রিডেক্স নিই। হয় $N_i$-তে
    তার অনুরূপ সমান কর্তনমাত্রার পদাংশ আছে, যা
    অনুমান অনুসারে $\cutrank{M}$-এর চেয়ে কম; নয়তো
    কর্তনমাত্রা $\len{!A_i}$, যা সংজ্ঞা অনুযায়ী
    $\cutrank{\ore{\ori{i}{}{\typeof{N}{!A_i}}}
      {\typeof{x_1}{!A_1}}{\typeof{N_1}{!C}}
      {\typeof{x_2}{!A_2}}{\typeof{N_2}{!C}}}$.
  \end{enumerate}
\end{proof}

\begin{thm}
  প্রত্যেক প্রমাণপদ স্বাভাবিক রূপে হ্রাস পায়;
  প্রত্যেক !!{derivation}s স্বাভাবিক !!{derivation}s-এ
  হ্রাস পায়।
\end{thm}

\begin{proof}
  প্রথমটি থেকে দ্বিতীয়টি পাওয়া যায়। $M$ প্রমাণপদ
  ধরে $m=\maxrank{M}$-এর উপর পূর্ণ আরোহে
  প্রথম দাবিটি প্রমাণ করি।
  \begin{enumerate}

  \item $m=0$ হলে $M$ ইতিমধ্যেই স্বাভাবিক।
  \item
    অন্যথায় $M$-এ কর্তনমাত্রা~$m$-যুক্ত
    রিডেক্সের সংখ্যা~$n$-এর উপর আরোহ করি।
    \begin{enumerate}
    \item $n=1$ হলে এমন রিডেক্স~$N$ বাছি,
      যাতে $m = \cutrank{N} > \cutrank{P}$,
      যখনই $P$ হলো~$N$-এর প্রকৃত রিডেক্স
      উপপদ। এমন পদাংশ থাকবেই, কারণ একটি পদে
      সসীমসংখ্যক উপপদ থাকে।

      $N$-এর সংকুচিত ফলকে $N'$ বলি। পূর্ববর্তী
      সহায়ক উপপাদ্য অনুযায়ী
      $\maxrank{N'} < \maxrank{N}$।
      তাই $\cutrank=m$-যুক্ত রিডেক্সের
      সংখ্যা~$n$ কমে। ফলে $m$-ও (অন্তত~$1$)
      কমে, এবং $m$-এর আরোহ-অনুমান প্রয়োগ করি।
    \item আরোহ-ধাপে $n>1$ ধরি। প্রক্রিয়াটি
      অনুরূপ, তবে $n$ কমে ধনাত্মক থাকে এবং
      তাই $m$ অপরিবর্তিত থাকে। তখন $n$-এর
      আরোহ-অনুমান প্রয়োগ করি।
    \end{enumerate}
\end{enumerate}
\end{proof}

টাইপযুক্ত $\lambd$-পদের স্বাভাবিকীকরণের অর্থ:
এমন প্রত্যেক পদের \emph{একটি স্বাভাবিক রূপ আছে}।
$\lambd$-পদে $\beta$-হ্রাস করে স্বাভাবিক রূপ
পাওয়াকে গণনা সম্পাদন এবং সেই রূপকে গণনার
মান ধরলে, টাইপযুক্ত $\lambd$-ক্যালকুলাসের
প্রত্যেক পদের জন্য অন্তত একটি গণনাপথ
শেষ পর্যন্ত থামে ও একটি মান দেয়।
তদুপরি টাইপ-সংরক্ষণ নিশ্চিত করে যে মানটির
টাইপ ঠিক আছে। যেমন, $N$-এর টাইপ
$!A \lif !B$ হলে (অর্থাৎ $!A$-টাইপের
যুক্তি নিয়ে $!B$-টাইপের মান দেওয়ার অপেক্ষক)
এবং $!A$-টাইপের পদ~$M$ (নিবেশ)-এ
তাকে প্রয়োগ করলে $(NM)$ হ্রাস পেয়ে
কোনো পদ~$N'$ (নির্গম)-এ পৌঁছায়; সেই
নির্গমের টাইপ যে~$!B$, তা নিশ্চিত।
% BN-SRC-611 TARGET-CORRECTION-BEGIN
\par\smallskip\noindent\textnormal{\small\textbf{উৎসসংশোধন (BN-SRC-611):} সদ্য প্রমাণিত স্বাভাবিক রূপের অস্তিত্ব থেকে কেবল অন্তত একটি সমাপ্ত হ্রাসপথ পাওয়া যায়; উৎসের “প্রত্যেক গণনা থামে” দাবিটি পরের শক্তিশালী স্বাভাবিকীকরণ ধর্মের বিষয়।}
% BN-SRC-611 TARGET-CORRECTION-END

বস্তুত টাইপযুক্ত $\lambd$-ক্যালকুলাসে উপপদে
$\beta$-হ্রাস প্রয়োগের শুধু একটি ক্রমেই
স্বাভাবিক রূপ পাওয়া যায়, এমন নয়; \emph{যেকোনো}
ক্রমে প্রয়োগ শেষ পর্যন্ত স্বাভাবিক রূপে থামে।
এই ধর্মকে \emph{শক্তিশালী স্বাভাবিকীকরণ}
বলে। অর্থাৎ গণনা থামিয়ে মান পাওয়ার
কেবল একটি নিশ্চিত পথই নেই; $\beta$-হ্রাসে গণনা চালানোর
\emph{প্রত্যেক} পথেই শেষ পর্যন্ত মান পাওয়া যায়।

গণনা চালানোর বিভিন্ন পথে \emph{একই}
মান পাওয়া যায় কীভাবে জানি? এখন পর্যন্ত
টাইপ-সংরক্ষণ থেকে শুধু জানি, পাওয়া সব
মানের টাইপ এক। টাইপযুক্ত $\lambd$-ক্যালকুলাসের
আরেকটি গুরুত্বপূর্ণ ধর্ম হলো \emph{সমাভিগম্যতা}
বা \emph{চার্চ--রসার ধর্ম}। $N \red N_1$
এবং $N \red N_2$ হলে এমন পদ~$N'$ আছে,
যাতে $N_1 \red N'$ এবং $N_2 \red N'$।
মনে করি, $N \red N'$ মানে শূন্য বা একাধিক
$\redone$-হ্রাসধাপে $N'$ পাওয়া যায়।
$N_1$ স্বাভাবিক রূপে থাকলে $N_1 \red N'$
পূরণকারী একমাত্র পদ~$N'$ হলো $N_1$ নিজেই
(শূন্য $\redone$-ধাপে)। ফলে $N_1$ ও
$N_2$ উভয়ই স্বাভাবিক রূপে থাকলে
$N_1 = N' = N_2$। অন্য কথায়,
$\redone$-এর শক্তিশালী স্বাভাবিকীকরণে
হ্রাসের প্রত্যেক পথ শেষ পর্যন্ত স্বাভাবিক রূপে
পৌঁছায়। $\red$-এর চার্চ--রসার ধর্মে
যেকোনো দুটি স্বাভাবিক রূপ সমান। অতএব
টাইপযুক্ত $\lambd$-ক্যালকুলাসে গণনা সব
পথেই থামে এবং তার মান (স্বাভাবিক রূপ)
গণনার পথের উপর নির্ভর করে না।
% BN-SRC-607 TARGET-CORRECTION-BEGIN
\par\smallskip\noindent\textnormal{\small\textbf{উৎসসংশোধন (BN-SRC-607):} শূন্য-ধাপ অনুমোদনকারী প্রতিফলিত-সঞ্চারী সম্পর্ক \(\red\)-এর বদলে শক্তিশালী স্বাভাবিকীকরণ এক-ধাপ সম্পর্ক \(\redone\)-এর ধর্ম হিসেবে নির্দিষ্ট করা হয়েছে।}
% BN-SRC-607 TARGET-CORRECTION-END

কোনো সম্পর্ক সমাভিগম্য প্রমাণ করা সাধারণত
কঠিন। তবে নিচের দুর্বলতর শর্তটি সহজেই
প্রতিষ্ঠা করতে পারি:

\begin{prop}[দুর্বল চার্চ--রসার ধর্ম]
$N \redone N_1$ এবং $N \redone N_2$
হলে এমন পদ~$N'$ আছে, যাতে $N_1 \red N'$
এবং $N_2 \red N'$।
\end{prop}

\begin{proof}
  $N$-এ যথাক্রমে রিডেক্স $M_1$ বা $M_2$
  সংকুচিত করে $N_1$ ও $N_2$ পাওয়া যায়।
  $M_1$ ও $M_2$ হলো $N$-এর উপপদ; তাদের
  অবস্থান হয় একটি অন্যটির ভিতরে, নয়তো
  পৃথক, অসংলগ্ন স্থানে। শেষের অবস্থাটি
  আগে দেখি: $N$-এর রূপ $N(M_1,M_2)$।
  $M_1$-কে তার সংকুচিত ফল~$M_1'$ দিয়ে
  বদলালে $N_1$, অর্থাৎ $N(M_1',M_2)$
  পাওয়া যায়; $M_2$-কে তার সংকুচিত ফল~$M_2'$
  দিয়ে বদলালে $N_2$, অর্থাৎ
  $N(M_1,M_2')$ পাওয়া যায়। দুটি পদই
  $N(M_1',M_2')$-এ হ্রাস পায়। অন্য
  অবস্থায় $M_2$-কে $M_1$-এর উপপদ ধরি
  (বিপরীত অবস্থাটি প্রতিসম); $M_1$ কোন
  ধরনের রিডেক্স, সেই অনুযায়ী
  অবস্থা আলাদা করি। যেমন,
  $M_1 \ident \proj{1}{\pair{O_1}{O_2}}$
  হলে তা $O_1$-এ হ্রাস পায়। $M_2$
  যদি $O_1$-এর উপপদ হয়, তবে $M_2$
  সংকুচিত করে $O_1'$ পাওয়া যায়। তাই
  আগে $M_1$, পরে $M_2$ সংকুচিত করলে
  ফল $O_1'$। অন্যদিকে আগে $M_2$
  সংকুচিত করলে
  $\proj{1}{\pair{O_1'}{O_2}}$
  পাওয়া যায়; এটিও $O_1'$-এ হ্রাস পায়।
% BN-SRC-608 TARGET-CORRECTION-BEGIN
\par\smallskip\noindent\textnormal{\small\textbf{উৎসসংশোধন (BN-SRC-608):} প্রথম প্রক্ষেপণের শেষ ফল উৎসে ভুল করে \(O_2\) ছিল; প্রদর্শিত পদটি \(O_1'\)-এ হ্রাস পায়।}
% BN-SRC-608 TARGET-CORRECTION-END
\end{proof}

\begin{lem}[নিউম্যানের সহায়ক উপপাদ্য]
$\redone$ শক্তিশালীভাবে স্বাভাবিকীকরণশীল হলে
এবং দুর্বল চার্চ--রসার ধর্ম থাকলে $\red$
সমাভিগম্য।
% BN-SRC-609 TARGET-CORRECTION-BEGIN
\par\smallskip\noindent\textnormal{\small\textbf{উৎসসংশোধন (BN-SRC-609):} শূন্য-ধাপসহ \(\red\) নয়, এক-ধাপ \(\redone\)-এর শক্তিশালী স্বাভাবিকীকরণ এখানে প্রয়োজন; সমাভিগম্যতা \(\red\)-এর। পরের প্রমাণে স্বাভাবিক রূপে পৌঁছানোর দাবি সর্বোচ্চ হ্রাস-অনুক্রমের জন্য।}
% BN-SRC-609 TARGET-CORRECTION-END
\end{lem}

\begin{proof}
  $\redone$ শক্তিশালীভাবে স্বাভাবিকীকরণশীল বলে
  প্রত্যেক সর্বোচ্চ হ্রাস-অনুক্রম স্বাভাবিক রূপে শেষ হয়।
  স্বাভাবিক রূপ অনন্য হওয়া এবং $\red$-এর
  সমাভিগম্যতা সমতুল্য। অতএব ধরি,
  কোনো পদ~$N$-এর দুটি ভিন্ন স্বাভাবিক রূপ
  $N'$ ও $N''$ আছে, নিচের অনুক্রমদুটির শেষে:
  \begin{align*}
  N & = N_1' \redone N_2' \redone \dots \redone N_k' = N' \text{ এবং}\\
  N & = N_1'' \redone N_2'' \redone \dots \redone N_l'' = N''
  \end{align*}
  (অনুক্রমদুটির দৈর্ঘ্য $>0$, নইলে $N$
  আগেই স্বাভাবিক রূপে থাকত।)

  $N_2' = N_2''$ হলে সেই প্রথম হ্রাসফলেরই
  দুটি ভিন্ন স্বাভাবিক রূপ $N'$ ও $N''$
  থাকে। $N_2' \neq N_2''$ হলে দুর্বল
  চার্চ--রসার ধর্মে এমন $N'''$ আছে,
  যাতে $N_2' \red N'''$ এবং
  $N_2'' \red N'''$। শক্তিশালী
  স্বাভাবিকীকরণে $N'''$-এর কোনো স্বাভাবিক
  রূপ~$Q$ আছে। $N' \neq N''$ বলে
  $Q \neq N'$ অথবা $Q \neq N''$।
  ফলে যথাক্রমে $N_2'$ অথবা $N_2''$-এর
  দুটি ভিন্ন স্বাভাবিক রূপ আছে। অর্থাৎ
  $N$-এর দুটি ভিন্ন স্বাভাবিক রূপ থাকলে
  এমন $N_1$ আছে, যাতে
  $N \redone N_1$ এবং $N_1$-এরও
  দুটি ভিন্ন স্বাভাবিক রূপ। একই যুক্তি
  চালিয়ে গেলে
  $N \redone N_1 \redone N_2 \redone \dots$
  অসীম অনুক্রম পাওয়া যায়; কিন্তু
  $\redone$ শক্তিশালীভাবে স্বাভাবিকীকরণশীল
  বলে তা অসম্ভব।
% BN-SRC-610 TARGET-CORRECTION-BEGIN
\par\smallskip\noindent\textnormal{\small\textbf{উৎসসংশোধন (BN-SRC-610):} দুই অনুক্রমের শুরু \(N_1'=N_1''=N\), তাই উৎসের সেই তুলনা ও তাদের দুর্বল-চার্চ--রসার প্রয়োগ অনর্থক। প্রথম হ্রাসফল \(N_2',N_2''\) তুলনা করা হয়েছে; তাদের সাধারণ পরবর্তী পদ \(N'''\)-এর একটি স্বাভাবিক রূপ \(Q\) নিয়ে ভিন্ন স্বাভাবিক রূপের বিরোধ সম্পন্ন করা হয়েছে।}
% BN-SRC-610 TARGET-CORRECTION-END
\end{proof}


\end{document}
